\label{sec:intervals-as-subsets}

\begin{topicbox}{Intervals as the Basic Pieces}
\begin{exposition}
The interval shapes were introduced in \S\ref{sec:intervals-on-a-number-line}. Here they are recast as subsets of the metric line and supplemented with boundedness, which is the size condition compactness will need.
\end{exposition}
\end{topicbox}
\begin{definitionbox}{Definition (Bounded Subset of the Real Line)}
\begin{definition}[Bounded Subset of the Real Line]
\label{def:bounded-subset-of-r}
A set $E\subseteq\mathbb{R}$ is \emph{bounded} exactly when there exists $M\ge 0$ such that $|x|\le M$ for all $x\in E$; equivalently, $E$ is contained in some closed interval $[-M,M]$.
\end{definition}
\end{definitionbox}
\begin{remark*}[Standard quantified statement]
\[
E \text{ is bounded} \Longleftrightarrow \exists M\ge 0\;\forall x\in E\;(|x|\le M).
\]
\end{remark*}
\begin{remark*}[Negated quantified statement]
\[
E \text{ is unbounded} \Longleftrightarrow \forall M\ge 0\;\exists x\in E\;(|x|>M).
\]
\end{remark*}
\begin{remark*}[Interpretation]
Boundedness says $E$ fits inside a single closed interval --- it does not run off to infinity. Together with closedness it is one half of the Heine--Borel characterization of compact subsets of $\mathbb{R}$.
\end{remark*}
\begin{dependencies}
\begin{itemize}
  \item \hyperref[def:closed-interval]{Closed Interval}
  \item \hyperref[thm:men-absolute-value-properties]{Absolute Value Properties}
\end{itemize}
\end{dependencies}
